% Section 1. Complementary Representations of Predictive Uncertainty
% Backbone: Discussion 6.1 (sec:discussion-representation). See STORY.md.

\section{Complementary Representations of Predictive Uncertainty}

% Slide 1.1. Source: Introduction. Background 2.6 (sec:entropy-decomposition),
% 2.7 (sec:bg-three-way-decomposition, eq:three-way-decomposition).
\begin{frame}{Uncertainty-Aware Classification}
	\textbf{Message.} A class label alone is not enough. We need the predictive uncertainty and its sources.

	\medskip
	\textbf{TODO: content.} Three uncertainty terms, one line each. Entropy decomposition in one line.
\end{frame}

% Slide 1.2. Source: Background 2.1 (sec:bg-bnn), 2.8
% (sec:bg-aleatoric-epistemic-decomposition, eq:point-mass-definition,
% eq:degenerate-distributional-term). Methods 4.1 (sec:bayesian-classification).
\begin{frame}{Categorical Bayesian Neural Network}
	\textbf{Message.} The Cat-BNN represents aleatoric and epistemic uncertainty. It represents no distributional uncertainty.

	\medskip
	\textbf{TODO: content.} Weight posterior, prediction $f(x;w)$, point mass on $\pi$, zero distributional term.
\end{frame}

% Slide 1.3. Source: Methods 4.2.1 (sec:edl-model-setup, eq:edl-dirichlet),
% 4.2.2 (sec:edl-representation).
\begin{frame}{Evidential Deep Learning}
	\textbf{Message.} EDL represents aleatoric and distributional uncertainty. It represents no epistemic uncertainty.

	\medskip
	\textbf{TODO: content.} Concentration parameters $\alpha(x;w)$, Dirichlet over $\pi$, one trained weight configuration.
\end{frame}

% Slide 1.4. Source: Introduction. Discussion 6.1. Methods 4.3 (sec:edl-bnn-setup),
% 4.4 (sec:count-likelihoods).
\begin{frame}{The Gap and Our Two Models}
	\textbf{Message.} We combine the weight posterior with the Dirichlet network. The result represents all three uncertainty terms.

	\medskip
	\textbf{TODO: content.} Table of models against uncertainty terms. Cat-eBNN for single labels. DM-eBNN for repeated labels.
\end{frame}
